\section{Connection to the Riemann Hypothesis}
\begin{quote}\small
\noindent\textbf{Status.} \emph{Legacy cross-domain bridge block inside the active main body.}\par
\noindent\textbf{Block function.} This section explains how the structural core is connected to the Riemann-side reading without turning the surrounding physical line into a separate proof source.
\end{quote}
\label{sec:rh}
%% ============================================================

\begin{remark}[Common origin of RH and entanglement]
The Riemann Hypothesis ($\zminus=0$ for all zeros) and quantum
entanglement (shared $\tau_n$) are both consequences of the same
biquaternion structure:
\begin{enumerate}
  \item The Minkowski norm $|q_n|^2=\tau_n^2$ is invariant
        (Theorem~\ref{thm:lorentz}).
  \item The cascade forces $\zminus\to 0$ along the light axis
        (Theorem~\ref{thm:fractal}).
  \item Stone--von Neumann makes all three kernels equivalent
        (Theorem~\ref{thm:unitary-equiv}).
\end{enumerate}
Neither was the goal of the original construction.
Both emerged from asking: what is the exact structure of the Sphaera?
\end{remark}

%% ============================================================

